%======================================================================
% _config_base.tex
% Configuration globale des documents CFPT
%======================================================================


%======================================================================
% 1 — TYPOGRAPHIE ET LANGUE
%======================================================================

\usepackage[T1]{fontenc}
\usepackage[french]{babel}
\usepackage[french]{varioref}

\usepackage{helvet}
\renewcommand{\familydefault}{\sfdefault}


%======================================================================
% 2 — MATHÉMATIQUES
%======================================================================

\usepackage{amsmath}
\usepackage{amsfonts}
\usepackage{amssymb}


%======================================================================
% 3 — GRAPHISMES
%======================================================================

\usepackage{graphicx}
\usepackage[table]{xcolor}
\usepackage{subcaption}

\DeclareGraphicsExtensions{.png,.PNG,.jpg,.JPG}

\graphicspath{{./medias/}{./captures/}{./includes/logo_default/}}


%======================================================================
% 4 — TABLEAUX ET MISE EN PAGE
%======================================================================

\usepackage{array}
\usepackage{tabularx}
\usepackage{multicol}
\usepackage{parskip}

\usepackage{geometry}

\geometry{
	left=2cm,
	right=2cm,
	top=2cm,
	bottom=2cm
}


%======================================================================
% 5 — BOÎTES / CODE / OUTILS
%======================================================================

\usepackage{fancyvrb}
\usepackage{marginnote}
\usepackage{tcolorbox}
\usepackage{tikz}
\usepackage{verbatim}
\usepackage{etoolbox}
\usepackage{minitoc}
\usepackage{float}
\usepackage{listings}

\tcbuselibrary{skins,breakable}


%======================================================================
% 6 — HYPERLIENS
%======================================================================

\usepackage[
colorlinks=true,
linkcolor=blue,
urlcolor=blue
]{hyperref}


%======================================================================
% 7 — DATE / HEURE
%======================================================================

\usepackage{datetime2}

\DTMsetdatestyle{ddmmyyyy}

\renewcommand{\DTMdisplaydate}[4]{#3/#2/#1}
\renewcommand{\DTMdisplaytime}[3]{#1:#2}

\DTMsetup{
	showzone=false,
	showseconds=false
}


%======================================================================
% 8 — CONFIGURATION DOCUMENTAIRE
%======================================================================

\newif\iftitlepagelogo

\usepackage{pgfkeys}

\pgfkeys{
	/doc/.is family,
	/doc/.cd,
	%------------------------------------------------
	% Toggles
	%------------------------------------------------
	titlepagelogo/.is if=titlepagelogo,
	%------------------------------------------------
	% Métadonnées
	%------------------------------------------------
	author/.estore in=\DocAuthor,
	email/.estore in=\DocEmail,
	title/.estore in=\DocTitle,
	subtitle/.estore in=\DocSubtitle,
	domaine/.estore in=\DocDomaine,
	year/.estore in=\DocYear,
	logo/.estore in=\DocLogo,
	keywords/.estore in=\DocKeywords,
	%------------------------------------------------
	% Configuration interne
	%------------------------------------------------
	embedsourcebasepath/.estore in=\embedsourcebasepath,
	%------------------------------------------------
	% HEADER / FOOTER CONFIGURABLE
	% Important : store in (pas estore in)
	% pour éviter l’expansion prématurée de \DTMnow,
	% \pageref, \rightmark, etc.
	%------------------------------------------------
	header-left/.store in=\DocHeaderLeft,
	header-center/.store in=\DocHeaderCenter,
	header-right/.store in=\DocHeaderRight,
	footer-left/.store in=\DocFooterLeft,
	footer-center/.store in=\DocFooterCenter,
	footer-right/.store in=\DocFooterRight
}

\newcommand{\DocSetup}[1]{\pgfkeys{/doc,#1}}


%======================================================================
% 9 — VALEURS PAR DÉFAUT
%======================================================================

\DocSetup{
	titlepagelogo=true,
	author={Stéphane Garchery},
	email={stephane.garchery@edu.ge.ch},
	title={Titre},
	subtitle={Sous titre},
	domaine={},
	year={2025--2026},
	logo={},
	keywords={},
	embedsourcebasepath={soluce},
	header-left={\footnotesize\textbf{\DocDomaine}},
	header-center={\footnotesize\DocTitle},
	header-right={\tiny\nouppercase{\rightmark}},
	footer-left={\tiny\DocEmail},
	footer-center={\tiny document généré le \DTMnow},
	footer-right={\scriptsize Page~\thepage~/~\pageref{LastPage}}
}


%======================================================================
% 10 — HEADERS / FOOTERS
%======================================================================

\usepackage{lastpage}
\usepackage{fancyhdr}

\pagestyle{fancy}

\fancyhf{}

\renewcommand{\headrulewidth}{0.4pt}
\renewcommand{\footrulewidth}{0.4pt}

\fancyhead[L]{\DocHeaderLeft}
\fancyhead[C]{\DocHeaderCenter}
\fancyhead[R]{\DocHeaderRight}

\fancyfoot[L]{\DocFooterLeft}
\fancyfoot[C]{\DocFooterCenter}
\fancyfoot[R]{\DocFooterRight}

\renewcommand{\sectionmark}[1]{\markright{\thesection\ #1}}


%======================================================================
% 11 — EXTENSIONS CFPT
%======================================================================

%======================================================================
% includes/inc_colors.tex
% Palette de couleurs centralisée pour tous les includes
%======================================================================

\usepackage{xcolor}

%----------------------------------------------------------------------
% Couleur de texte "atténué" (muted)
%----------------------------------------------------------------------
\definecolor{mutedColor}{gray}{0.45} % texte secondaire (pied de page, légendes…)

%----------------------------------------------------------------------
% Gris / neutres de base
%----------------------------------------------------------------------
\definecolor{cfptBlack}{rgb}{0.00,0.00,0.00}
\definecolor{cfptDarkGray}{gray}{0.25}
\definecolor{cfptGray}{gray}{0.50}
\definecolor{cfptLightGray}{gray}{0.90}
\definecolor{cfptVeryLightGray}{gray}{0.98}

% Alias "historiques" pour compatibilité
\colorlet{lightgray}{cfptVeryLightGray}
\colorlet{gray}{cfptGray}

%======================================================================
% Palette pour les listings de code
%======================================================================

% Fond des blocs de code
\colorlet{codeBackground}{cfptVeryLightGray}

% Couleurs sémantiques pour la coloration syntaxique
\definecolor{codeKeyword}{rgb}{0.00,0.00,0.80}    % bleu foncé
\definecolor{codeIdentifier}{rgb}{0.00,0.60,0.00} % vert
\definecolor{codeString}{rgb}{0.80,0.00,0.00}     % rouge
\definecolor{codeComment}{rgb}{0.50,0.50,0.50}    % gris
\definecolor{codeAccent}{rgb}{0.58,0.00,0.82}     % violet (true/false/null)
\definecolor{codeOperator}{rgb}{0.90,0.45,0.00}   % orange (opérateurs, ponctuation)
\definecolor{codeType}{rgb}{0.20,0.60,0.60}       % cyan pour types/classes

% Alias "historiques" pour compatibilité avec les anciens includes
\colorlet{dkblue}{codeKeyword}
\colorlet{dkgreen}{codeIdentifier}
\colorlet{red}{codeString}
\colorlet{purple}{codeAccent}
\colorlet{orange}{codeOperator}   % !!! réservé à la coloration de code
\colorlet{class}{codeType}
\colorlet{darkgreen}{codeIdentifier}
\colorlet{darkred}{codeString}

%======================================================================
% Palette pour les boîtes (info / success / warning / danger / etc.)
%  -> Inspirée de Bootstrap (blue / green / amber / red) avec fonds doux.
%======================================================================

% Couleurs pour la boîte par défaut
\definecolor{defaultBackground}{RGB}{245,245,245} 		% gris très clair
\colorlet{defaultColor}{gray}                     		% bordure grise
\colorlet{defaultTitleBackground}{gray!15}        		% fond du titre légèrement grisé

% INFO  : messages neutres / informatifs
\definecolor{infoColor}{rgb}{0.00,0.45,0.80}            % bleu
\definecolor{infoBackground}{rgb}{0.91,0.96,1.00}       % bleu très clair
\definecolor{infoTitleBackground}{rgb}{0.00,0.45,0.80}  % bandeau titre

% SUCCESS : succès, validation, confirmation
\definecolor{successColor}{rgb}{0.17,0.63,0.34}               % vert
\definecolor{successBackground}{rgb}{0.91,0.98,0.92}          % vert très clair
\definecolor{successTitleBackground}{rgb}{0.17,0.63,0.34}

% WARNING : avertissement, point de vigilance
\definecolor{warningColor}{rgb}{0.98,0.76,0.18}               % jaune / amber
\definecolor{warningBackground}{rgb}{1.00,0.97,0.88}          % jaune très clair
\definecolor{warningTitleBackground}{rgb}{0.98,0.76,0.18}

% DANGER : erreur critique, interdiction
\definecolor{dangerColor}{rgb}{0.83,0.25,0.25}                % rouge
\definecolor{dangerBackground}{rgb}{1.00,0.93,0.94}           % rouge très clair
\definecolor{dangerTitleBackground}{rgb}{0.83,0.25,0.25}

% TEACHER : notes internes formateur
\definecolor{teacherColor}{rgb}{0.55,0.27,0.07}         	% brun soutenu
\definecolor{teacherBackground}{rgb}{0.98,0.95,0.90}    	% beige très clair
\definecolor{teacherTitleBackground}{rgb}{0.55,0.27,0.07}

% DOWNLOAD : encadrés de téléchargement / ressources
\definecolor{downloadColor}{rgb}{0.00,0.40,0.80}              % bleu soutenu
\definecolor{downloadBackground}{rgb}{0.92,0.96,1.00}         % bleu clair
\definecolor{downloadTitleBackground}{rgb}{0.00,0.40,0.80}

% AIDE : indices, coups de pouce
\definecolor{aideColor}{rgb}{0.40,0.23,0.72}                  % indigo / violet
\definecolor{aideBackground}{rgb}{0.95,0.93,1.00}             % indigo très clair
\definecolor{aideTitleBackground}{rgb}{0.40,0.23,0.72}

% OBJECTIF : objectifs pédagogiques, buts de séance
\definecolor{objectifColor}{rgb}{1.00,0.58,0.00}              % orange franc
\definecolor{objectifBackground}{rgb}{1.00,0.95,0.90}         % orange très clair
\definecolor{objectifTitleBackground}{rgb}{1.00,0.58,0.00}

% VALIDATION : critères de validation, checklist finale
\definecolor{validationColor}{rgb}{0.00,0.58,0.51}            % vert/bleu (teal)
\definecolor{validationBackground}{rgb}{0.90,0.97,0.96}       % teal très clair
\definecolor{validationTitleBackground}{rgb}{0.00,0.58,0.51}

% (Pas d’alias ici sur 'orange' ou 'cyan' pour éviter de casser les listings)

%======================================================================
% includes/inc_listings.tex
% Configuration centralisée pour tous les listings de code
%
% DÉPENDANCES :
%   - nécessite que les couleurs soient définies dans includes/inc_colors.tex
%     (par exemple : dkblue, dkgreen, darkred, class, gray, lightgray, etc.)
%
% CHARGEMENT TYPE DANS LE DOCUMENT :
%   %======================================================================
% includes/inc_colors.tex
% Palette de couleurs centralisée pour tous les includes
%======================================================================

\usepackage{xcolor}

%----------------------------------------------------------------------
% Couleur de texte "atténué" (muted)
%----------------------------------------------------------------------
\definecolor{mutedColor}{gray}{0.45} % texte secondaire (pied de page, légendes…)

%----------------------------------------------------------------------
% Gris / neutres de base
%----------------------------------------------------------------------
\definecolor{cfptBlack}{rgb}{0.00,0.00,0.00}
\definecolor{cfptDarkGray}{gray}{0.25}
\definecolor{cfptGray}{gray}{0.50}
\definecolor{cfptLightGray}{gray}{0.90}
\definecolor{cfptVeryLightGray}{gray}{0.98}

% Alias "historiques" pour compatibilité
\colorlet{lightgray}{cfptVeryLightGray}
\colorlet{gray}{cfptGray}

%======================================================================
% Palette pour les listings de code
%======================================================================

% Fond des blocs de code
\colorlet{codeBackground}{cfptVeryLightGray}

% Couleurs sémantiques pour la coloration syntaxique
\definecolor{codeKeyword}{rgb}{0.00,0.00,0.80}    % bleu foncé
\definecolor{codeIdentifier}{rgb}{0.00,0.60,0.00} % vert
\definecolor{codeString}{rgb}{0.80,0.00,0.00}     % rouge
\definecolor{codeComment}{rgb}{0.50,0.50,0.50}    % gris
\definecolor{codeAccent}{rgb}{0.58,0.00,0.82}     % violet (true/false/null)
\definecolor{codeOperator}{rgb}{0.90,0.45,0.00}   % orange (opérateurs, ponctuation)
\definecolor{codeType}{rgb}{0.20,0.60,0.60}       % cyan pour types/classes

% Alias "historiques" pour compatibilité avec les anciens includes
\colorlet{dkblue}{codeKeyword}
\colorlet{dkgreen}{codeIdentifier}
\colorlet{red}{codeString}
\colorlet{purple}{codeAccent}
\colorlet{orange}{codeOperator}   % !!! réservé à la coloration de code
\colorlet{class}{codeType}
\colorlet{darkgreen}{codeIdentifier}
\colorlet{darkred}{codeString}

%======================================================================
% Palette pour les boîtes (info / success / warning / danger / etc.)
%  -> Inspirée de Bootstrap (blue / green / amber / red) avec fonds doux.
%======================================================================

% Couleurs pour la boîte par défaut
\definecolor{defaultBackground}{RGB}{245,245,245} 		% gris très clair
\colorlet{defaultColor}{gray}                     		% bordure grise
\colorlet{defaultTitleBackground}{gray!15}        		% fond du titre légèrement grisé

% INFO  : messages neutres / informatifs
\definecolor{infoColor}{rgb}{0.00,0.45,0.80}            % bleu
\definecolor{infoBackground}{rgb}{0.91,0.96,1.00}       % bleu très clair
\definecolor{infoTitleBackground}{rgb}{0.00,0.45,0.80}  % bandeau titre

% SUCCESS : succès, validation, confirmation
\definecolor{successColor}{rgb}{0.17,0.63,0.34}               % vert
\definecolor{successBackground}{rgb}{0.91,0.98,0.92}          % vert très clair
\definecolor{successTitleBackground}{rgb}{0.17,0.63,0.34}

% WARNING : avertissement, point de vigilance
\definecolor{warningColor}{rgb}{0.98,0.76,0.18}               % jaune / amber
\definecolor{warningBackground}{rgb}{1.00,0.97,0.88}          % jaune très clair
\definecolor{warningTitleBackground}{rgb}{0.98,0.76,0.18}

% DANGER : erreur critique, interdiction
\definecolor{dangerColor}{rgb}{0.83,0.25,0.25}                % rouge
\definecolor{dangerBackground}{rgb}{1.00,0.93,0.94}           % rouge très clair
\definecolor{dangerTitleBackground}{rgb}{0.83,0.25,0.25}

% TEACHER : notes internes formateur
\definecolor{teacherColor}{rgb}{0.55,0.27,0.07}         	% brun soutenu
\definecolor{teacherBackground}{rgb}{0.98,0.95,0.90}    	% beige très clair
\definecolor{teacherTitleBackground}{rgb}{0.55,0.27,0.07}

% DOWNLOAD : encadrés de téléchargement / ressources
\definecolor{downloadColor}{rgb}{0.00,0.40,0.80}              % bleu soutenu
\definecolor{downloadBackground}{rgb}{0.92,0.96,1.00}         % bleu clair
\definecolor{downloadTitleBackground}{rgb}{0.00,0.40,0.80}

% AIDE : indices, coups de pouce
\definecolor{aideColor}{rgb}{0.40,0.23,0.72}                  % indigo / violet
\definecolor{aideBackground}{rgb}{0.95,0.93,1.00}             % indigo très clair
\definecolor{aideTitleBackground}{rgb}{0.40,0.23,0.72}

% OBJECTIF : objectifs pédagogiques, buts de séance
\definecolor{objectifColor}{rgb}{1.00,0.58,0.00}              % orange franc
\definecolor{objectifBackground}{rgb}{1.00,0.95,0.90}         % orange très clair
\definecolor{objectifTitleBackground}{rgb}{1.00,0.58,0.00}

% VALIDATION : critères de validation, checklist finale
\definecolor{validationColor}{rgb}{0.00,0.58,0.51}            % vert/bleu (teal)
\definecolor{validationBackground}{rgb}{0.90,0.97,0.96}       % teal très clair
\definecolor{validationTitleBackground}{rgb}{0.00,0.58,0.51}

% (Pas d’alias ici sur 'orange' ou 'cyan' pour éviter de casser les listings)

%   %======================================================================
% includes/inc_listings.tex
% Configuration centralisée pour tous les listings de code
%
% DÉPENDANCES :
%   - nécessite que les couleurs soient définies dans includes/inc_colors.tex
%     (par exemple : dkblue, dkgreen, darkred, class, gray, lightgray, etc.)
%
% CHARGEMENT TYPE DANS LE DOCUMENT :
%   %======================================================================
% includes/inc_colors.tex
% Palette de couleurs centralisée pour tous les includes
%======================================================================

\usepackage{xcolor}

%----------------------------------------------------------------------
% Couleur de texte "atténué" (muted)
%----------------------------------------------------------------------
\definecolor{mutedColor}{gray}{0.45} % texte secondaire (pied de page, légendes…)

%----------------------------------------------------------------------
% Gris / neutres de base
%----------------------------------------------------------------------
\definecolor{cfptBlack}{rgb}{0.00,0.00,0.00}
\definecolor{cfptDarkGray}{gray}{0.25}
\definecolor{cfptGray}{gray}{0.50}
\definecolor{cfptLightGray}{gray}{0.90}
\definecolor{cfptVeryLightGray}{gray}{0.98}

% Alias "historiques" pour compatibilité
\colorlet{lightgray}{cfptVeryLightGray}
\colorlet{gray}{cfptGray}

%======================================================================
% Palette pour les listings de code
%======================================================================

% Fond des blocs de code
\colorlet{codeBackground}{cfptVeryLightGray}

% Couleurs sémantiques pour la coloration syntaxique
\definecolor{codeKeyword}{rgb}{0.00,0.00,0.80}    % bleu foncé
\definecolor{codeIdentifier}{rgb}{0.00,0.60,0.00} % vert
\definecolor{codeString}{rgb}{0.80,0.00,0.00}     % rouge
\definecolor{codeComment}{rgb}{0.50,0.50,0.50}    % gris
\definecolor{codeAccent}{rgb}{0.58,0.00,0.82}     % violet (true/false/null)
\definecolor{codeOperator}{rgb}{0.90,0.45,0.00}   % orange (opérateurs, ponctuation)
\definecolor{codeType}{rgb}{0.20,0.60,0.60}       % cyan pour types/classes

% Alias "historiques" pour compatibilité avec les anciens includes
\colorlet{dkblue}{codeKeyword}
\colorlet{dkgreen}{codeIdentifier}
\colorlet{red}{codeString}
\colorlet{purple}{codeAccent}
\colorlet{orange}{codeOperator}   % !!! réservé à la coloration de code
\colorlet{class}{codeType}
\colorlet{darkgreen}{codeIdentifier}
\colorlet{darkred}{codeString}

%======================================================================
% Palette pour les boîtes (info / success / warning / danger / etc.)
%  -> Inspirée de Bootstrap (blue / green / amber / red) avec fonds doux.
%======================================================================

% Couleurs pour la boîte par défaut
\definecolor{defaultBackground}{RGB}{245,245,245} 		% gris très clair
\colorlet{defaultColor}{gray}                     		% bordure grise
\colorlet{defaultTitleBackground}{gray!15}        		% fond du titre légèrement grisé

% INFO  : messages neutres / informatifs
\definecolor{infoColor}{rgb}{0.00,0.45,0.80}            % bleu
\definecolor{infoBackground}{rgb}{0.91,0.96,1.00}       % bleu très clair
\definecolor{infoTitleBackground}{rgb}{0.00,0.45,0.80}  % bandeau titre

% SUCCESS : succès, validation, confirmation
\definecolor{successColor}{rgb}{0.17,0.63,0.34}               % vert
\definecolor{successBackground}{rgb}{0.91,0.98,0.92}          % vert très clair
\definecolor{successTitleBackground}{rgb}{0.17,0.63,0.34}

% WARNING : avertissement, point de vigilance
\definecolor{warningColor}{rgb}{0.98,0.76,0.18}               % jaune / amber
\definecolor{warningBackground}{rgb}{1.00,0.97,0.88}          % jaune très clair
\definecolor{warningTitleBackground}{rgb}{0.98,0.76,0.18}

% DANGER : erreur critique, interdiction
\definecolor{dangerColor}{rgb}{0.83,0.25,0.25}                % rouge
\definecolor{dangerBackground}{rgb}{1.00,0.93,0.94}           % rouge très clair
\definecolor{dangerTitleBackground}{rgb}{0.83,0.25,0.25}

% TEACHER : notes internes formateur
\definecolor{teacherColor}{rgb}{0.55,0.27,0.07}         	% brun soutenu
\definecolor{teacherBackground}{rgb}{0.98,0.95,0.90}    	% beige très clair
\definecolor{teacherTitleBackground}{rgb}{0.55,0.27,0.07}

% DOWNLOAD : encadrés de téléchargement / ressources
\definecolor{downloadColor}{rgb}{0.00,0.40,0.80}              % bleu soutenu
\definecolor{downloadBackground}{rgb}{0.92,0.96,1.00}         % bleu clair
\definecolor{downloadTitleBackground}{rgb}{0.00,0.40,0.80}

% AIDE : indices, coups de pouce
\definecolor{aideColor}{rgb}{0.40,0.23,0.72}                  % indigo / violet
\definecolor{aideBackground}{rgb}{0.95,0.93,1.00}             % indigo très clair
\definecolor{aideTitleBackground}{rgb}{0.40,0.23,0.72}

% OBJECTIF : objectifs pédagogiques, buts de séance
\definecolor{objectifColor}{rgb}{1.00,0.58,0.00}              % orange franc
\definecolor{objectifBackground}{rgb}{1.00,0.95,0.90}         % orange très clair
\definecolor{objectifTitleBackground}{rgb}{1.00,0.58,0.00}

% VALIDATION : critères de validation, checklist finale
\definecolor{validationColor}{rgb}{0.00,0.58,0.51}            % vert/bleu (teal)
\definecolor{validationBackground}{rgb}{0.90,0.97,0.96}       % teal très clair
\definecolor{validationTitleBackground}{rgb}{0.00,0.58,0.51}

% (Pas d’alias ici sur 'orange' ou 'cyan' pour éviter de casser les listings)

%   %======================================================================
% includes/inc_listings.tex
% Configuration centralisée pour tous les listings de code
%
% DÉPENDANCES :
%   - nécessite que les couleurs soient définies dans includes/inc_colors.tex
%     (par exemple : dkblue, dkgreen, darkred, class, gray, lightgray, etc.)
%
% CHARGEMENT TYPE DANS LE DOCUMENT :
%   \input{includes/inc_colors}
%   \input{includes/inc_listings}
%
% USAGE GÉNÉRAL :
%   - Style par défaut neutre :
%       \begin{lstlisting}
%       ...
%       \end{lstlisting}
%
%   - C# (style csharp) :
%       \begin{lstlisting}[style=csharp]
%       ...
%       \end{lstlisting}
%
%   - PHP (style php) :
%       \begin{lstlisting}[style=php]
%       ...
%       \end{lstlisting}
%
%   - JavaScript (style javascript) :
%       \begin{lstlisting}[style=javascript]
%       ...
%       \end{lstlisting}
%
%   - SQL / MySQL (style sql) :
%       \begin{lstlisting}[style=sql]
%       ...
%       \end{lstlisting}
%
%   - Bash (style bashsimple) :
%       \begin{lstlisting}[style=bashsimple]
%       ...
%   \end{lstlisting}
%
%   - Pseudo-code (style pseudocode) :
%       \begin{lstlisting}[style=pseudocode]
%       ...
%   \end{lstlisting}
%
%   - Listing protégé contre le copier/coller :
%       \begin{lstlistingnocopy}{Texte-remplacant-dans-le-presse-papiers}{CSharp}
%       ...
%   \end{lstlistingnocopy}
%
%======================================================================

%----------------------------------------------------------------------
% Tentative de chargement automatique de inc_colors si nécessaire
% On teste une couleur définie dans inc_colors (par ex. dkblue).
%----------------------------------------------------------------------
\makeatletter
\@ifundefined{color@dkblue}{%
	% dkblue non défini -> on suppose que inc_colors n'a pas encore été chargé
	\InputIfFileExists{includes/inc_colors.tex}{}{%
		% Si le fichier n'existe pas, on ne fait rien (fallback silencieux)
	}%
}{}%
\makeatother

%----------------------------------------------------------------------
% Packages nécessaires
%----------------------------------------------------------------------
\usepackage{listings}
\usepackage{listingsutf8} % charge listings + support UTF-8
\usepackage{accsupp}      % pour masquer le texte lors du copy/paste (lstlistingnocopy)

%======================================================================
% PARTIE 1 : Configuration commune
%======================================================================

% Numéros de ligne non copiables (pour certains styles, ex. C#)
\newcommand{\noncopynumber}[1]{%
	\BeginAccSupp{method=escape,ActualText={}}%
	{\scriptsize #1}%
	\EndAccSupp{}%
}

% Réglages globaux pour tous les listings
\lstset{
	basicstyle       = \footnotesize\ttfamily,
	columns          = fullflexible,
	keepspaces       = true,
	showstringspaces = false,
	upquote          = true,
	breaklines       = true,
	breakatwhitespace= true,
	tabsize          = 2,
	frame            = single,
	captionpos       = b,
	inputencoding    = utf8,
	extendedchars    = true,
	% Cartographie commune (accents FR + ponctuation typographique)
	literate         =
		{á}{{\'a}}1 {à}{{\`a}}1 {â}{{\^a}}1 {ä}{{\"a}}1
		{Á}{{\'A}}1 {À}{{\`A}}1 {Â}{{\^A}}1 {Ä}{{\"A}}1
		{é}{{\'e}}1 {è}{{\`e}}1 {ê}{{\^e}}1 {ë}{{\"e}}1
		{É}{{\'E}}1 {È}{{\`E}}1 {Ê}{{\^E}}1 {Ë}{{\"E}}1
		{í}{{\'i}}1 {ì}{{\`i}}1 {î}{{\^i}}1 {ï}{{\"i}}1
		{Í}{{\'I}}1 {Ì}{{\`I}}1 {Î}{{\^I}}1 {Ï}{{\"I}}1
		{ó}{{\'o}}1 {ò}{{\`o}}1 {ô}{{\^o}}1 {ö}{{\"o}}1
		{Ó}{{\'O}}1 {Ò}{{\`O}}1 {Ô}{{\^O}}1 {Ö}{{\"O}}1
		{ú}{{\'u}}1 {ù}{{\`u}}1 {û}{{\^u}}1 {ü}{{\"u}}1
		{Ú}{{\'U}}1 {Ù}{{\`U}}1 {Û}{{\^U}}1 {Ü}{{\"U}}1
		{ç}{{\c{c}}}1 {Ç}{{\c{C}}}1
		{œ}{{\oe}}1 {Œ}{{\OE}}1
		{æ}{{\ae}}1 {Æ}{{\AE}}1
		{’}{{'}}1 {‘}{{`}}1
		{“}{{``}}1 {”}{{''}}1
		{–}{{--}}1 {—}{{---}}1
		{…}{{\ldots}}1
	,
	backgroundcolor  = \color{lightgray},
	rulesepcolor     = \color{gray},
	prebreak         = \raisebox{0ex}[0ex][0ex]{\ensuremath{\hookleftarrow}},
	% Inclusion partielle via rangeprefix/rangesuffix
	rangeprefix        = {//`:~},
	rangesuffix        = {~:`},
	includerangemarker = false,
}

% Style par défaut neutre
\lstdefinestyle{defaultcode}{
	keywordstyle    = \color{dkblue}\bfseries,
	identifierstyle = \color{dkgreen},
	stringstyle     = \color{red},
	commentstyle    = \color{gray}\itshape,
}

\lstset{style=defaultcode}

%======================================================================
% PARTIE 2 : C# (CSharp)
%======================================================================

% Langage C# basé sur [Sharp]C
\lstdefinelanguage{CSharp}{
	language      = [Sharp]C,
	morekeywords  = {
		if, for, while, do, var, get, set,
		string, value, String, ToString, ToInt32, ToBool, toInt
	},
	otherkeywords = {
		\#region, \#endregion, \#define, \#if, \#endif, \#else
	},
}

% Style CSharp complet
\lstdefinestyle{csharp}{
	language        = CSharp,
	numbers         = left,
	numberstyle     = \noncopynumber,
	stepnumber      = 5,
	numbersep       = 5pt,
	frame           = shadowbox,
	keywordstyle    = \color{dkblue}\bfseries,
	commentstyle    = \color{darkgreen},
	stringstyle     = \color{darkred},
	identifierstyle = \color{black},
	% Classes/types mis en avant
	emph            = [1]{
		Application,
		Char, Color, Console, Convert,
		DialogResult,
		Environment, EventArgs,
		Form,
		Object, String,
		SByte, Int16, Int32, Int64,
		Byte, UInt16, UInt32, UInt64,
		Single, Double,
		Keys, KeyPressEventArgs,
		MessageBox, MessageBoxButtons,
		% --- MSTest : classes / attributs importants ---
		Assert,
		TestContext,
		TestClass, TestMethod,
		TestInitialize, TestCleanup,
		ClassInitialize, ClassCleanup,
		DataTestMethod, DataRow,
		ExpectedException,
		TestCategory, Owner, Priority, TestProperty,
		% --- MonoGame / XNA : classes / structs / enums importants ---
		Game, GameTime, GameComponent, DrawableGameComponent,
		GraphicsDevice, GraphicsDeviceManager, GraphicsProfile,
		Viewport, DisplayMode, PresentationParameters,
		ContentManager, ContentReader, ContentWriter,
		SpriteBatch, SpriteFont, SpriteEffects, SpriteSortMode,
		Texture, Texture2D, Texture3D, RenderTarget2D, RenderTargetCube,
		SurfaceFormat, RenderTargetUsage,
		Effect, BasicEffect, AlphaTestEffect,
		BlendState, SamplerState, DepthStencilState, RasterizerState,
		Color, Rectangle, Point,
		Vector2, Vector3, Vector4, Matrix, Quaternion, Plane, Ray,
		MathHelper, Curve, BoundingBox, BoundingSphere, BoundingFrustum,
		GameServiceContainer,
		% --- Input ---
		Keyboard, KeyboardState,
		Mouse, MouseState,
		GamePad, GamePadState, GamePadCapabilities,
		Keys, Buttons, ButtonState, PlayerIndex,
		TouchPanel, TouchCollection, TouchLocation, GestureSample, GestureType,
		% --- Audio / Media ---
		SoundEffect, SoundEffectInstance, SoundState,
		Song, MediaPlayer,
		% --- Timing / misc ---
		TimeSpan,
		% --- Common graphics helpers (si tu utilises les primitives) ---
		VertexDeclaration,
		VertexBuffer, IndexBuffer, DynamicVertexBuffer,
		VertexPositionColor, VertexPositionTexture, VertexPositionColorTexture,
		PrimitiveType, BufferUsage, IndexElementSize
	},
	emphstyle       = [1]\color{class},
}


%======================================================================
% PARTIE 3 : PHP
%======================================================================

% Charger la définition standard de PHP
\lstloadlanguages{PHP}

% Sous-langage PHP "CFPT" : on ajoute nos mots-clés / types
\lstdefinelanguage[CFPT]{PHP}[]{PHP}{%
	% Mots-clés supplémentaires
	morekeywords={declare},
	% Types scalaires / retour qu'on veut colorer comme "types"
	morekeywords=[2]{%
		int, float, string, bool, boolean, void,
		array, iterable, object, mixed, never
	},
}

\lstdefinestyle{php}{%
	language        = [CFPT]PHP,
	% groupe 1 : mots-clés PHP
	keywordstyle    = \color{dkblue}\bfseries,
	% groupe 2 : types
	keywordstyle    = [2]\color{class}\bfseries,
	identifierstyle = \color{dkgreen},
	stringstyle     = \color{red},
	commentstyle    = \color{gray}\itshape,
	emph            = {true, false, null},
	emphstyle       = \color{purple},
	morekeywords    = {
		namespace, use, class, extends,
		public, private, protected, function,
		return, as, new, static, global, var,
		if, else, elseif, while, for, foreach, do,
		switch, case, break, continue, default,
		echo, print, die, exit, instanceof,
		try, catch, finally, throw, implements
	},
	% Délimiteurs spéciaux @@...@@ (par ex. pour marquer du code important)
	moredelim       = **[is][\color{orange}]{@@}{@@},
}

%======================================================================
% PARTIE 4 : JavaScript
%======================================================================

% Définition du langage JavaScript pour listings
\lstdefinelanguage{JavaScript}{
	sensitive        = true,
	% Mots-clés principaux JavaScript
	morekeywords = {
		break, case, catch, continue, debugger, default, delete, do, else, finally,
		for, function, if, in, instanceof, new, return, switch, this, throw, try,
		typeof, var, void, while, with,
		const, let,
		class, extends, super, constructor,
		import, export, from, as,
		get, set,
		async, await, yield
	},
	% Types / "mots-clés" de type (pratique si tu montres un peu de TS ou JSDoc)
	morekeywords = [2]{
		number, string, boolean, any, void, unknown, never, object, Promise
	},
	% Mots-clés / fonctions spécifiques à Jest (groupe 3)
	morekeywords = [3]{
		describe, it, test,
		expect,
		beforeEach, afterEach,
		beforeAll, afterAll,
		jest
	},
	% Commentaires
	morecomment = [l]{//},
	morecomment = [s]{/*}{*/},
	% Chaînes
	morestring  = [b]',
	morestring  = [b]"
}

% Style JavaScript (couleurs cohérentes avec le reste)
\lstdefinestyle{javascript}{%
	language        = JavaScript,
	basicstyle      = \footnotesize\ttfamily,
	% groupe 1 : mots-clés JS
	keywordstyle    = \color{dkblue}\bfseries,
	% groupe 2 : types / notions proches TypeScript
	keywordstyle    = [2]\color{class}\bfseries,
	% groupe 3 : Jest (API de tests)
	keywordstyle    = [3]\color{darkred}\bfseries,
	identifierstyle = \color{dkgreen},
	stringstyle     = \color{red},
	commentstyle    = \color{gray}\itshape,
	emph            = {true, false, null, undefined},
	emphstyle       = \color{purple},
}

%======================================================================
% PARTIE 5 : SQL / MySQL
%======================================================================

\lstdefinelanguage{mysql}{
	morekeywords = {
		% --- requêtes de base ---
		SELECT, FROM, WHERE, INSERT, INTO, VALUES, UPDATE, DELETE,
		CREATE, TABLE, PRIMARY, KEY, AUTO_INCREMENT,
		ALTER, DROP, DATABASE, USE,
		INNER, LEFT, RIGHT, JOIN, ON, AS,
		AND, OR, NOT, NULL, IS, IN, LIKE, BETWEEN,
		ORDER, BY, GROUP, HAVING, DISTINCT, LIMIT, OFFSET,
		FOR,
		% --- transactions / verrous ---
		BEGIN, START, TRANSACTION, COMMIT, ROLLBACK,
		SAVEPOINT, RELEASE, WORK, CHAIN,
		LOCK, UNLOCK,
		% --- niveaux d'isolation / modes d'accès ---
		ISOLATION, LEVEL,
		READ, WRITE, ONLY,
		REPEATABLE, UNCOMMITTED, COMMITTED, SERIALIZABLE,
		AUTOCOMMIT,
		% --- XA transactions ---
		XA, END, PREPARE, RECOVER,
		% --- types ---
		INT, INTEGER, DECIMAL, NUMERIC, FLOAT, DOUBLE,
		CHAR, VARCHAR, TEXT, DATE, DATETIME, TIMESTAMP, TIME,
		% --- contraintes / index ---
		ENGINE, DEFAULT, UNIQUE, INDEX, FOREIGN, REFERENCES,
		% --- gestion des comptes / privilèges ---
		GRANT, REVOKE, ALL, PRIVILEGES,
		USER, ROLE, PASSWORD,
		IDENTIFIED, BY,
		FLUSH, WITH, OPTION, PROXY,
		% --- vues ---
		VIEW, VIEWS,
		% --- procédures / fonctions / triggers / events ---
		PROCEDURE, FUNCTION, TRIGGER, EVENT, ROUTINE,
		INOUT, RETURNS, DETERMINISTIC,
		SQL, SECURITY, DEFINER, INVOKER,
		SET
	},
	sensitive   = false, % mots-clés insensibles à la casse
}

\lstdefinestyle{sql}{
	language        = mysql,
	keywordstyle    = \color{dkblue}\bfseries,
	identifierstyle = \color{dkgreen},
	stringstyle     = \color{red},
	commentstyle    = \color{gray}\itshape,
}

%======================================================================
% PARTIE 6 : Bash / PowerShell (style commun)
%======================================================================

\lstdefinestyle{bashsimple}{
	language        = bash,
	basicstyle      = \footnotesize\ttfamily,
	frame           = single,
	framesep        = 2pt,
	breaklines      = true,
	showstringspaces= false,
	columns         = fullflexible,
	keepspaces      = true,
	upquote         = true,
	tabsize         = 2,
	keywordstyle    = \color{dkblue}\bfseries,
	identifierstyle = \color{black},
	stringstyle     = \color{red},
	commentstyle    = \color{gray}\itshape,
	% IMPORTANT : le tiret fait partie des "lettres"
	alsoletter      = {-},
	% force en noir litéralement
	% IMPORTANT : on conserve la cartographie FR + on ajoute les règles spécifiques
	literate=
		{á}{{\'a}}1 {à}{{\`a}}1 {â}{{\^a}}1 {ä}{{\"a}}1
		{Á}{{\'A}}1 {À}{{\`A}}1 {Â}{{\^A}}1 {Ä}{{\"A}}1
		{é}{{\'e}}1 {è}{{\`e}}1 {ê}{{\^e}}1 {ë}{{\"e}}1
		{É}{{\'E}}1 {È}{{\`E}}1 {Ê}{{\^E}}1 {Ë}{{\"E}}1
		{í}{{\'i}}1 {ì}{{\`i}}1 {î}{{\^i}}1 {ï}{{\"i}}1
		{Í}{{\'I}}1 {Ì}{{\`I}}1 {Î}{{\^I}}1 {Ï}{{\"I}}1
		{ó}{{\'o}}1 {ò}{{\`o}}1 {ô}{{\^o}}1 {ö}{{\"o}}1
		{Ó}{{\'O}}1 {Ò}{{\`O}}1 {Ô}{{\^O}}1 {Ö}{{\"O}}1
		{ú}{{\'u}}1 {ù}{{\`u}}1 {û}{{\^u}}1 {ü}{{\"u}}1
		{Ú}{{\'U}}1 {Ù}{{\`U}}1 {Û}{{\^U}}1 {Ü}{{\"U}}1
		{ç}{{\c{c}}}1 {Ç}{{\c{C}}}1
		{/php/}{{{\color{black}/php/}}}4%
		{/php*}{{{\color{black}/php/}}}4%
		{php.ini}{{{\color{black}php.ini}}}6%
		{.php}{{{\color{black}.php}}}4%
		{wsl.}{{{\color{black}wsl.}}}4
	,
	%--------------------------------------------------------------
	% Commandes bash / sh + cmdlets PowerShell mises en évidence
	% (j'ai retiré 'tar' pour éviter de colorer les extensions .tar)
	%--------------------------------------------------------------
	morekeywords={
		% --- Bash / sh de base ---
		clear,
		cd, ls, cp, mv, rm, mkdir, rmdir, touch,
		chmod, chown, chgrp, umask, usermod,
		sudo, su, passwd,
		apt, apt-get, apt-cache, add-apt-repository,
		a2enmod, a2dismod, a2encond, a2dismod, a2enconf, a2disconf,
		service, systemctl,
		kill, killall, ps, top, htop,
		cat, tac, less, more, head, tail,
		grep, egrep, fgrep, find, xargs,
		echo, printf,
		nano, vi, vim, code,
		ssh, scp, sftp, rsync,
		service,
		zip, unzip, gzip, gunzip,
		df, du, mount, umount,
		uname, hostname, whoami, id, pwd,
		ln,
		php, php8.4,
		mysql, mariadb, mariadb_secure_installation, mysql_secure_installation,
		update-alternatives,
		composer,
		wsl,
		% --- Cmdlets PowerShell fréquents ---
		Get-Help, Get-Command,
		Get-ChildItem, Get-Item, Set-Item, Remove-Item,
		Copy-Item, Move-Item, New-Item, Test-Path,
		Get-Location, Set-Location, Push-Location, Pop-Location,
		Get-Content, Set-Content, Add-Content,
		Get-Process, Stop-Process, Start-Process,
		Get-Service, Start-Service, Stop-Service, Restart-Service,
		Get-Module, Import-Module, Install-Module,
		Invoke-WebRequest, Invoke-Expression,
		Get-History, Clear-History,
		Get-Alias, Set-Alias,
		Get-ExecutionPolicy, Set-ExecutionPolicy,
		Get-EventLog, Write-EventLog
	},
	%--------------------------------------------------------------
	% Options / paramètres mis en avant (bash + PowerShell)
	%--------------------------------------------------------------
	emph={
		% options bash/génériques
		-d,
		-s, status,
		-c,--config,
		-h,--help,
		-v,--version,
		-y,--yes,
		-r,-R,
		-a,--all,
		-l,--list,
		-f,--force,
		--update,
		-i,--ini,
		--install,install,
		--export, --terminate,
		--import,
		-name,
		restart,
		--unregister,
		--set-default,
		--list,
		--verbose,
		-G,
		-u, --user,
		-p, --password,
		update, full-upgrade,
		% options courantes PowerShell
		-Path,
		-Recurse,
		-Filter,
		-Include,
		-Exclude,
		-Force,
		-Name,
		-ErrorAction,
		-WhatIf,
		-Confirm,
		-File,
		-Directory,
		-Value,
		-Destination,
		-Source,
		-Credential
	},
	emphstyle=\color{purple}\bfseries,
}

%======================================================================
% Style listings : Pseudo-code (Flowgorithm / ATALGO)
%======================================================================

\lstdefinestyle{pseudocode}{
	language        = {},
	basicstyle      = \footnotesize\ttfamily,
	showstringspaces=false,
	columns         = fullflexible,
	keepspaces      = true,
	upquote         = true,
	% Commentaires (si tu en mets dans tes extraits)
	morecomment=[l]{//},
	morecomment=[s]{/*}{*/},
	commentstyle    = \color{gray}\itshape,
	% Chaînes
	morestring=[b]",
	stringstyle     = \color{darkred},
	% -------------------------
	% Groupe 1 : Structure / contrôle
	% -------------------------
	keywordstyle    = \color{dkblue}\bfseries,
	morekeywords    = {
		ALGORITHME, DEBUT_ALGORITHME, FIN_ALGORITHME,
		Fonction, Principale, DEBUT, FIN, Fin,
		SI, Si, ALORS, Alors, SINON, Sinon, Autre, DEBUT_SI, FIN_SI, DEBUT_SINON, FIN_SINON,
		TANT_QUE, Tant, Que, DEBUT_TANT_QUE, FIN_TANT_QUE,
		POUR, Pour, DEBUT_POUR, FIN_POUR, ALLANT_DE, PAR_PAS_DE, A, a,
		FAIRE, Faire,
		CAS, SUIVANT, DEBUT_SUIVANT, FIN_SUIVANT
	},
	% -------------------------
	% Groupe 2 : I/O et actions “verbes” (couleur différente)
	% -------------------------
	classoffset=1,
	keywordstyle=\color{darkgreen}\bfseries,
	morekeywords={
		AFFICHER, Afficher,
		LIRE, Lire,
		ATTENDRE, Attendre,
		Saisie, SAISIE,
		Appeler, APPELER,
		Initialiser, INITIALISER
	},
	classoffset=0,
	% -------------------------
	% Groupe 3 : Déclarations / affectations (encore une autre couleur)
	% -------------------------
	classoffset=2,
	keywordstyle=\color{dkblue}\bfseries,
	morekeywords={
		VARIABLES, Variables,
		Déclarer, Declarer, DECLARER,
		Variable,
		Type, TYPE, de,
		Assigner, ASSIGNER,
		Expression, EXPRESSION,
		Paramètre, Parametre, PARAMETRE,
		Retour, RETOUR,
		Retourner, Retourner, RETOURNER,
		Retouner, RETOUNER, % (typo présent dans le PDF)
		Prend, PREND, PREND_POUR_VALEUR
	},
	classoffset=0,
	% -------------------------
	% Types fréquents (mis en “emphase”)
	% -------------------------
	emph={Entier,Réel,Reel,Chaîne,Chaine,Booléen,Booleen,Tableau,Array},
	emphstyle=\color{darkred}\bfseries,
	% Fonctions usuelles vues dans tes pseudo-codes
	emph={size,len,char,Random,True,False},
	emphstyle=\color{darkgreen},
	% Bonus : rendre quelques opérateurs plus “propres” (optionnel)
	literate=
		{á}{{\'a}}1 {à}{{\`a}}1 {â}{{\^a}}1 {ä}{{\"a}}1
		{Á}{{\'A}}1 {À}{{\`A}}1 {Â}{{\^A}}1 {Ä}{{\"A}}1
		{é}{{\'e}}1 {è}{{\`e}}1 {ê}{{\^e}}1 {ë}{{\"e}}1
		{É}{{\'E}}1 {È}{{\`E}}1 {Ê}{{\^E}}1 {Ë}{{\"E}}1
		{í}{{\'i}}1 {ì}{{\`i}}1 {î}{{\^i}}1 {ï}{{\"i}}1
		{Í}{{\'I}}1 {Ì}{{\`I}}1 {Î}{{\^I}}1 {Ï}{{\"I}}1
		{ó}{{\'o}}1 {ò}{{\`o}}1 {ô}{{\^o}}1 {ö}{{\"o}}1
		{Ó}{{\'O}}1 {Ò}{{\`O}}1 {Ô}{{\^O}}1 {Ö}{{\"O}}1
		{ú}{{\'u}}1 {ù}{{\`u}}1 {û}{{\^u}}1 {ü}{{\"u}}1
		{Ú}{{\'U}}1 {Ù}{{\`U}}1 {Û}{{\^U}}1 {Ü}{{\"U}}1
		{ç}{{\c{c}}}1 {Ç}{{\c{C}}}1
		{<=}{{\(\le\)}}1
		{>=}{{\(\ge\)}}1
		{!=}{{\(\neq\)}}1
		{&&}{{\(\wedge\)}}1
		{||}{{\(\vee\)}}1
}

%=========================================================
% Langage LaTeX
%=========================================================
\lstdefinelanguage{latex}{
	morekeywords={
		documentclass,usepackage,input,include,includegraphics,
		begin,end,section,subsection,subsubsection,
		tableofcontents,listoffigures,lstlistoflistings,
		caption,label,ref,vref,textbf,textit,emph,
		DocSetup,newcommand,renewcommand,
		graphicspath,DeclareGraphicsExtensions
	},
	sensitive=true,
	morecomment=[l]{\%},
	morestring=[b]",
	alsoletter={\\}
}

%=========================================================
% Style LaTeX
%=========================================================
\lstdefinestyle{latexdoc}{
	language=latex,
	basicstyle=\ttfamily\small,
	keywordstyle=\color{blue}\bfseries,
	commentstyle=\color{gray},
	stringstyle=\color{red},
	numbers=left,
	numberstyle=\tiny,
	stepnumber=1,
	numbersep=8pt,
	frame=single,
	breaklines=true,
	showstringspaces=false,
	tabsize=2
}

%======================================================================
% PARTIE 8 : Environnement protégé (no copy/paste)
%======================================================================

\lstnewenvironment{lstlistingnocopy}[2]{%
	% #1 : texte remplacé dans le presse-papiers
	% #2 : langage ou style (ex. CSharp, php, mysql, ...)
	\BeginAccSupp{method=escape,ActualText={#1}}%
	\lstset{
		language = #2,
	}%
}{%
	\EndAccSupp{}%
}

%
% USAGE GÉNÉRAL :
%   - Style par défaut neutre :
%       \begin{lstlisting}
%       ...
%       \end{lstlisting}
%
%   - C# (style csharp) :
%       \begin{lstlisting}[style=csharp]
%       ...
%       \end{lstlisting}
%
%   - PHP (style php) :
%       \begin{lstlisting}[style=php]
%       ...
%       \end{lstlisting}
%
%   - JavaScript (style javascript) :
%       \begin{lstlisting}[style=javascript]
%       ...
%       \end{lstlisting}
%
%   - SQL / MySQL (style sql) :
%       \begin{lstlisting}[style=sql]
%       ...
%       \end{lstlisting}
%
%   - Bash (style bashsimple) :
%       \begin{lstlisting}[style=bashsimple]
%       ...
%   \end{lstlisting}
%
%   - Pseudo-code (style pseudocode) :
%       \begin{lstlisting}[style=pseudocode]
%       ...
%   \end{lstlisting}
%
%   - Listing protégé contre le copier/coller :
%       \begin{lstlistingnocopy}{Texte-remplacant-dans-le-presse-papiers}{CSharp}
%       ...
%   \end{lstlistingnocopy}
%
%======================================================================

%----------------------------------------------------------------------
% Tentative de chargement automatique de inc_colors si nécessaire
% On teste une couleur définie dans inc_colors (par ex. dkblue).
%----------------------------------------------------------------------
\makeatletter
\@ifundefined{color@dkblue}{%
	% dkblue non défini -> on suppose que inc_colors n'a pas encore été chargé
	\InputIfFileExists{includes/inc_colors.tex}{}{%
		% Si le fichier n'existe pas, on ne fait rien (fallback silencieux)
	}%
}{}%
\makeatother

%----------------------------------------------------------------------
% Packages nécessaires
%----------------------------------------------------------------------
\usepackage{listings}
\usepackage{listingsutf8} % charge listings + support UTF-8
\usepackage{accsupp}      % pour masquer le texte lors du copy/paste (lstlistingnocopy)

%======================================================================
% PARTIE 1 : Configuration commune
%======================================================================

% Numéros de ligne non copiables (pour certains styles, ex. C#)
\newcommand{\noncopynumber}[1]{%
	\BeginAccSupp{method=escape,ActualText={}}%
	{\scriptsize #1}%
	\EndAccSupp{}%
}

% Réglages globaux pour tous les listings
\lstset{
	basicstyle       = \footnotesize\ttfamily,
	columns          = fullflexible,
	keepspaces       = true,
	showstringspaces = false,
	upquote          = true,
	breaklines       = true,
	breakatwhitespace= true,
	tabsize          = 2,
	frame            = single,
	captionpos       = b,
	inputencoding    = utf8,
	extendedchars    = true,
	% Cartographie commune (accents FR + ponctuation typographique)
	literate         =
		{á}{{\'a}}1 {à}{{\`a}}1 {â}{{\^a}}1 {ä}{{\"a}}1
		{Á}{{\'A}}1 {À}{{\`A}}1 {Â}{{\^A}}1 {Ä}{{\"A}}1
		{é}{{\'e}}1 {è}{{\`e}}1 {ê}{{\^e}}1 {ë}{{\"e}}1
		{É}{{\'E}}1 {È}{{\`E}}1 {Ê}{{\^E}}1 {Ë}{{\"E}}1
		{í}{{\'i}}1 {ì}{{\`i}}1 {î}{{\^i}}1 {ï}{{\"i}}1
		{Í}{{\'I}}1 {Ì}{{\`I}}1 {Î}{{\^I}}1 {Ï}{{\"I}}1
		{ó}{{\'o}}1 {ò}{{\`o}}1 {ô}{{\^o}}1 {ö}{{\"o}}1
		{Ó}{{\'O}}1 {Ò}{{\`O}}1 {Ô}{{\^O}}1 {Ö}{{\"O}}1
		{ú}{{\'u}}1 {ù}{{\`u}}1 {û}{{\^u}}1 {ü}{{\"u}}1
		{Ú}{{\'U}}1 {Ù}{{\`U}}1 {Û}{{\^U}}1 {Ü}{{\"U}}1
		{ç}{{\c{c}}}1 {Ç}{{\c{C}}}1
		{œ}{{\oe}}1 {Œ}{{\OE}}1
		{æ}{{\ae}}1 {Æ}{{\AE}}1
		{’}{{'}}1 {‘}{{`}}1
		{“}{{``}}1 {”}{{''}}1
		{–}{{--}}1 {—}{{---}}1
		{…}{{\ldots}}1
	,
	backgroundcolor  = \color{lightgray},
	rulesepcolor     = \color{gray},
	prebreak         = \raisebox{0ex}[0ex][0ex]{\ensuremath{\hookleftarrow}},
	% Inclusion partielle via rangeprefix/rangesuffix
	rangeprefix        = {//`:~},
	rangesuffix        = {~:`},
	includerangemarker = false,
}

% Style par défaut neutre
\lstdefinestyle{defaultcode}{
	keywordstyle    = \color{dkblue}\bfseries,
	identifierstyle = \color{dkgreen},
	stringstyle     = \color{red},
	commentstyle    = \color{gray}\itshape,
}

\lstset{style=defaultcode}

%======================================================================
% PARTIE 2 : C# (CSharp)
%======================================================================

% Langage C# basé sur [Sharp]C
\lstdefinelanguage{CSharp}{
	language      = [Sharp]C,
	morekeywords  = {
		if, for, while, do, var, get, set,
		string, value, String, ToString, ToInt32, ToBool, toInt
	},
	otherkeywords = {
		\#region, \#endregion, \#define, \#if, \#endif, \#else
	},
}

% Style CSharp complet
\lstdefinestyle{csharp}{
	language        = CSharp,
	numbers         = left,
	numberstyle     = \noncopynumber,
	stepnumber      = 5,
	numbersep       = 5pt,
	frame           = shadowbox,
	keywordstyle    = \color{dkblue}\bfseries,
	commentstyle    = \color{darkgreen},
	stringstyle     = \color{darkred},
	identifierstyle = \color{black},
	% Classes/types mis en avant
	emph            = [1]{
		Application,
		Char, Color, Console, Convert,
		DialogResult,
		Environment, EventArgs,
		Form,
		Object, String,
		SByte, Int16, Int32, Int64,
		Byte, UInt16, UInt32, UInt64,
		Single, Double,
		Keys, KeyPressEventArgs,
		MessageBox, MessageBoxButtons,
		% --- MSTest : classes / attributs importants ---
		Assert,
		TestContext,
		TestClass, TestMethod,
		TestInitialize, TestCleanup,
		ClassInitialize, ClassCleanup,
		DataTestMethod, DataRow,
		ExpectedException,
		TestCategory, Owner, Priority, TestProperty,
		% --- MonoGame / XNA : classes / structs / enums importants ---
		Game, GameTime, GameComponent, DrawableGameComponent,
		GraphicsDevice, GraphicsDeviceManager, GraphicsProfile,
		Viewport, DisplayMode, PresentationParameters,
		ContentManager, ContentReader, ContentWriter,
		SpriteBatch, SpriteFont, SpriteEffects, SpriteSortMode,
		Texture, Texture2D, Texture3D, RenderTarget2D, RenderTargetCube,
		SurfaceFormat, RenderTargetUsage,
		Effect, BasicEffect, AlphaTestEffect,
		BlendState, SamplerState, DepthStencilState, RasterizerState,
		Color, Rectangle, Point,
		Vector2, Vector3, Vector4, Matrix, Quaternion, Plane, Ray,
		MathHelper, Curve, BoundingBox, BoundingSphere, BoundingFrustum,
		GameServiceContainer,
		% --- Input ---
		Keyboard, KeyboardState,
		Mouse, MouseState,
		GamePad, GamePadState, GamePadCapabilities,
		Keys, Buttons, ButtonState, PlayerIndex,
		TouchPanel, TouchCollection, TouchLocation, GestureSample, GestureType,
		% --- Audio / Media ---
		SoundEffect, SoundEffectInstance, SoundState,
		Song, MediaPlayer,
		% --- Timing / misc ---
		TimeSpan,
		% --- Common graphics helpers (si tu utilises les primitives) ---
		VertexDeclaration,
		VertexBuffer, IndexBuffer, DynamicVertexBuffer,
		VertexPositionColor, VertexPositionTexture, VertexPositionColorTexture,
		PrimitiveType, BufferUsage, IndexElementSize
	},
	emphstyle       = [1]\color{class},
}


%======================================================================
% PARTIE 3 : PHP
%======================================================================

% Charger la définition standard de PHP
\lstloadlanguages{PHP}

% Sous-langage PHP "CFPT" : on ajoute nos mots-clés / types
\lstdefinelanguage[CFPT]{PHP}[]{PHP}{%
	% Mots-clés supplémentaires
	morekeywords={declare},
	% Types scalaires / retour qu'on veut colorer comme "types"
	morekeywords=[2]{%
		int, float, string, bool, boolean, void,
		array, iterable, object, mixed, never
	},
}

\lstdefinestyle{php}{%
	language        = [CFPT]PHP,
	% groupe 1 : mots-clés PHP
	keywordstyle    = \color{dkblue}\bfseries,
	% groupe 2 : types
	keywordstyle    = [2]\color{class}\bfseries,
	identifierstyle = \color{dkgreen},
	stringstyle     = \color{red},
	commentstyle    = \color{gray}\itshape,
	emph            = {true, false, null},
	emphstyle       = \color{purple},
	morekeywords    = {
		namespace, use, class, extends,
		public, private, protected, function,
		return, as, new, static, global, var,
		if, else, elseif, while, for, foreach, do,
		switch, case, break, continue, default,
		echo, print, die, exit, instanceof,
		try, catch, finally, throw, implements
	},
	% Délimiteurs spéciaux @@...@@ (par ex. pour marquer du code important)
	moredelim       = **[is][\color{orange}]{@@}{@@},
}

%======================================================================
% PARTIE 4 : JavaScript
%======================================================================

% Définition du langage JavaScript pour listings
\lstdefinelanguage{JavaScript}{
	sensitive        = true,
	% Mots-clés principaux JavaScript
	morekeywords = {
		break, case, catch, continue, debugger, default, delete, do, else, finally,
		for, function, if, in, instanceof, new, return, switch, this, throw, try,
		typeof, var, void, while, with,
		const, let,
		class, extends, super, constructor,
		import, export, from, as,
		get, set,
		async, await, yield
	},
	% Types / "mots-clés" de type (pratique si tu montres un peu de TS ou JSDoc)
	morekeywords = [2]{
		number, string, boolean, any, void, unknown, never, object, Promise
	},
	% Mots-clés / fonctions spécifiques à Jest (groupe 3)
	morekeywords = [3]{
		describe, it, test,
		expect,
		beforeEach, afterEach,
		beforeAll, afterAll,
		jest
	},
	% Commentaires
	morecomment = [l]{//},
	morecomment = [s]{/*}{*/},
	% Chaînes
	morestring  = [b]',
	morestring  = [b]"
}

% Style JavaScript (couleurs cohérentes avec le reste)
\lstdefinestyle{javascript}{%
	language        = JavaScript,
	basicstyle      = \footnotesize\ttfamily,
	% groupe 1 : mots-clés JS
	keywordstyle    = \color{dkblue}\bfseries,
	% groupe 2 : types / notions proches TypeScript
	keywordstyle    = [2]\color{class}\bfseries,
	% groupe 3 : Jest (API de tests)
	keywordstyle    = [3]\color{darkred}\bfseries,
	identifierstyle = \color{dkgreen},
	stringstyle     = \color{red},
	commentstyle    = \color{gray}\itshape,
	emph            = {true, false, null, undefined},
	emphstyle       = \color{purple},
}

%======================================================================
% PARTIE 5 : SQL / MySQL
%======================================================================

\lstdefinelanguage{mysql}{
	morekeywords = {
		% --- requêtes de base ---
		SELECT, FROM, WHERE, INSERT, INTO, VALUES, UPDATE, DELETE,
		CREATE, TABLE, PRIMARY, KEY, AUTO_INCREMENT,
		ALTER, DROP, DATABASE, USE,
		INNER, LEFT, RIGHT, JOIN, ON, AS,
		AND, OR, NOT, NULL, IS, IN, LIKE, BETWEEN,
		ORDER, BY, GROUP, HAVING, DISTINCT, LIMIT, OFFSET,
		FOR,
		% --- transactions / verrous ---
		BEGIN, START, TRANSACTION, COMMIT, ROLLBACK,
		SAVEPOINT, RELEASE, WORK, CHAIN,
		LOCK, UNLOCK,
		% --- niveaux d'isolation / modes d'accès ---
		ISOLATION, LEVEL,
		READ, WRITE, ONLY,
		REPEATABLE, UNCOMMITTED, COMMITTED, SERIALIZABLE,
		AUTOCOMMIT,
		% --- XA transactions ---
		XA, END, PREPARE, RECOVER,
		% --- types ---
		INT, INTEGER, DECIMAL, NUMERIC, FLOAT, DOUBLE,
		CHAR, VARCHAR, TEXT, DATE, DATETIME, TIMESTAMP, TIME,
		% --- contraintes / index ---
		ENGINE, DEFAULT, UNIQUE, INDEX, FOREIGN, REFERENCES,
		% --- gestion des comptes / privilèges ---
		GRANT, REVOKE, ALL, PRIVILEGES,
		USER, ROLE, PASSWORD,
		IDENTIFIED, BY,
		FLUSH, WITH, OPTION, PROXY,
		% --- vues ---
		VIEW, VIEWS,
		% --- procédures / fonctions / triggers / events ---
		PROCEDURE, FUNCTION, TRIGGER, EVENT, ROUTINE,
		INOUT, RETURNS, DETERMINISTIC,
		SQL, SECURITY, DEFINER, INVOKER,
		SET
	},
	sensitive   = false, % mots-clés insensibles à la casse
}

\lstdefinestyle{sql}{
	language        = mysql,
	keywordstyle    = \color{dkblue}\bfseries,
	identifierstyle = \color{dkgreen},
	stringstyle     = \color{red},
	commentstyle    = \color{gray}\itshape,
}

%======================================================================
% PARTIE 6 : Bash / PowerShell (style commun)
%======================================================================

\lstdefinestyle{bashsimple}{
	language        = bash,
	basicstyle      = \footnotesize\ttfamily,
	frame           = single,
	framesep        = 2pt,
	breaklines      = true,
	showstringspaces= false,
	columns         = fullflexible,
	keepspaces      = true,
	upquote         = true,
	tabsize         = 2,
	keywordstyle    = \color{dkblue}\bfseries,
	identifierstyle = \color{black},
	stringstyle     = \color{red},
	commentstyle    = \color{gray}\itshape,
	% IMPORTANT : le tiret fait partie des "lettres"
	alsoletter      = {-},
	% force en noir litéralement
	% IMPORTANT : on conserve la cartographie FR + on ajoute les règles spécifiques
	literate=
		{á}{{\'a}}1 {à}{{\`a}}1 {â}{{\^a}}1 {ä}{{\"a}}1
		{Á}{{\'A}}1 {À}{{\`A}}1 {Â}{{\^A}}1 {Ä}{{\"A}}1
		{é}{{\'e}}1 {è}{{\`e}}1 {ê}{{\^e}}1 {ë}{{\"e}}1
		{É}{{\'E}}1 {È}{{\`E}}1 {Ê}{{\^E}}1 {Ë}{{\"E}}1
		{í}{{\'i}}1 {ì}{{\`i}}1 {î}{{\^i}}1 {ï}{{\"i}}1
		{Í}{{\'I}}1 {Ì}{{\`I}}1 {Î}{{\^I}}1 {Ï}{{\"I}}1
		{ó}{{\'o}}1 {ò}{{\`o}}1 {ô}{{\^o}}1 {ö}{{\"o}}1
		{Ó}{{\'O}}1 {Ò}{{\`O}}1 {Ô}{{\^O}}1 {Ö}{{\"O}}1
		{ú}{{\'u}}1 {ù}{{\`u}}1 {û}{{\^u}}1 {ü}{{\"u}}1
		{Ú}{{\'U}}1 {Ù}{{\`U}}1 {Û}{{\^U}}1 {Ü}{{\"U}}1
		{ç}{{\c{c}}}1 {Ç}{{\c{C}}}1
		{/php/}{{{\color{black}/php/}}}4%
		{/php*}{{{\color{black}/php/}}}4%
		{php.ini}{{{\color{black}php.ini}}}6%
		{.php}{{{\color{black}.php}}}4%
		{wsl.}{{{\color{black}wsl.}}}4
	,
	%--------------------------------------------------------------
	% Commandes bash / sh + cmdlets PowerShell mises en évidence
	% (j'ai retiré 'tar' pour éviter de colorer les extensions .tar)
	%--------------------------------------------------------------
	morekeywords={
		% --- Bash / sh de base ---
		clear,
		cd, ls, cp, mv, rm, mkdir, rmdir, touch,
		chmod, chown, chgrp, umask, usermod,
		sudo, su, passwd,
		apt, apt-get, apt-cache, add-apt-repository,
		a2enmod, a2dismod, a2encond, a2dismod, a2enconf, a2disconf,
		service, systemctl,
		kill, killall, ps, top, htop,
		cat, tac, less, more, head, tail,
		grep, egrep, fgrep, find, xargs,
		echo, printf,
		nano, vi, vim, code,
		ssh, scp, sftp, rsync,
		service,
		zip, unzip, gzip, gunzip,
		df, du, mount, umount,
		uname, hostname, whoami, id, pwd,
		ln,
		php, php8.4,
		mysql, mariadb, mariadb_secure_installation, mysql_secure_installation,
		update-alternatives,
		composer,
		wsl,
		% --- Cmdlets PowerShell fréquents ---
		Get-Help, Get-Command,
		Get-ChildItem, Get-Item, Set-Item, Remove-Item,
		Copy-Item, Move-Item, New-Item, Test-Path,
		Get-Location, Set-Location, Push-Location, Pop-Location,
		Get-Content, Set-Content, Add-Content,
		Get-Process, Stop-Process, Start-Process,
		Get-Service, Start-Service, Stop-Service, Restart-Service,
		Get-Module, Import-Module, Install-Module,
		Invoke-WebRequest, Invoke-Expression,
		Get-History, Clear-History,
		Get-Alias, Set-Alias,
		Get-ExecutionPolicy, Set-ExecutionPolicy,
		Get-EventLog, Write-EventLog
	},
	%--------------------------------------------------------------
	% Options / paramètres mis en avant (bash + PowerShell)
	%--------------------------------------------------------------
	emph={
		% options bash/génériques
		-d,
		-s, status,
		-c,--config,
		-h,--help,
		-v,--version,
		-y,--yes,
		-r,-R,
		-a,--all,
		-l,--list,
		-f,--force,
		--update,
		-i,--ini,
		--install,install,
		--export, --terminate,
		--import,
		-name,
		restart,
		--unregister,
		--set-default,
		--list,
		--verbose,
		-G,
		-u, --user,
		-p, --password,
		update, full-upgrade,
		% options courantes PowerShell
		-Path,
		-Recurse,
		-Filter,
		-Include,
		-Exclude,
		-Force,
		-Name,
		-ErrorAction,
		-WhatIf,
		-Confirm,
		-File,
		-Directory,
		-Value,
		-Destination,
		-Source,
		-Credential
	},
	emphstyle=\color{purple}\bfseries,
}

%======================================================================
% Style listings : Pseudo-code (Flowgorithm / ATALGO)
%======================================================================

\lstdefinestyle{pseudocode}{
	language        = {},
	basicstyle      = \footnotesize\ttfamily,
	showstringspaces=false,
	columns         = fullflexible,
	keepspaces      = true,
	upquote         = true,
	% Commentaires (si tu en mets dans tes extraits)
	morecomment=[l]{//},
	morecomment=[s]{/*}{*/},
	commentstyle    = \color{gray}\itshape,
	% Chaînes
	morestring=[b]",
	stringstyle     = \color{darkred},
	% -------------------------
	% Groupe 1 : Structure / contrôle
	% -------------------------
	keywordstyle    = \color{dkblue}\bfseries,
	morekeywords    = {
		ALGORITHME, DEBUT_ALGORITHME, FIN_ALGORITHME,
		Fonction, Principale, DEBUT, FIN, Fin,
		SI, Si, ALORS, Alors, SINON, Sinon, Autre, DEBUT_SI, FIN_SI, DEBUT_SINON, FIN_SINON,
		TANT_QUE, Tant, Que, DEBUT_TANT_QUE, FIN_TANT_QUE,
		POUR, Pour, DEBUT_POUR, FIN_POUR, ALLANT_DE, PAR_PAS_DE, A, a,
		FAIRE, Faire,
		CAS, SUIVANT, DEBUT_SUIVANT, FIN_SUIVANT
	},
	% -------------------------
	% Groupe 2 : I/O et actions “verbes” (couleur différente)
	% -------------------------
	classoffset=1,
	keywordstyle=\color{darkgreen}\bfseries,
	morekeywords={
		AFFICHER, Afficher,
		LIRE, Lire,
		ATTENDRE, Attendre,
		Saisie, SAISIE,
		Appeler, APPELER,
		Initialiser, INITIALISER
	},
	classoffset=0,
	% -------------------------
	% Groupe 3 : Déclarations / affectations (encore une autre couleur)
	% -------------------------
	classoffset=2,
	keywordstyle=\color{dkblue}\bfseries,
	morekeywords={
		VARIABLES, Variables,
		Déclarer, Declarer, DECLARER,
		Variable,
		Type, TYPE, de,
		Assigner, ASSIGNER,
		Expression, EXPRESSION,
		Paramètre, Parametre, PARAMETRE,
		Retour, RETOUR,
		Retourner, Retourner, RETOURNER,
		Retouner, RETOUNER, % (typo présent dans le PDF)
		Prend, PREND, PREND_POUR_VALEUR
	},
	classoffset=0,
	% -------------------------
	% Types fréquents (mis en “emphase”)
	% -------------------------
	emph={Entier,Réel,Reel,Chaîne,Chaine,Booléen,Booleen,Tableau,Array},
	emphstyle=\color{darkred}\bfseries,
	% Fonctions usuelles vues dans tes pseudo-codes
	emph={size,len,char,Random,True,False},
	emphstyle=\color{darkgreen},
	% Bonus : rendre quelques opérateurs plus “propres” (optionnel)
	literate=
		{á}{{\'a}}1 {à}{{\`a}}1 {â}{{\^a}}1 {ä}{{\"a}}1
		{Á}{{\'A}}1 {À}{{\`A}}1 {Â}{{\^A}}1 {Ä}{{\"A}}1
		{é}{{\'e}}1 {è}{{\`e}}1 {ê}{{\^e}}1 {ë}{{\"e}}1
		{É}{{\'E}}1 {È}{{\`E}}1 {Ê}{{\^E}}1 {Ë}{{\"E}}1
		{í}{{\'i}}1 {ì}{{\`i}}1 {î}{{\^i}}1 {ï}{{\"i}}1
		{Í}{{\'I}}1 {Ì}{{\`I}}1 {Î}{{\^I}}1 {Ï}{{\"I}}1
		{ó}{{\'o}}1 {ò}{{\`o}}1 {ô}{{\^o}}1 {ö}{{\"o}}1
		{Ó}{{\'O}}1 {Ò}{{\`O}}1 {Ô}{{\^O}}1 {Ö}{{\"O}}1
		{ú}{{\'u}}1 {ù}{{\`u}}1 {û}{{\^u}}1 {ü}{{\"u}}1
		{Ú}{{\'U}}1 {Ù}{{\`U}}1 {Û}{{\^U}}1 {Ü}{{\"U}}1
		{ç}{{\c{c}}}1 {Ç}{{\c{C}}}1
		{<=}{{\(\le\)}}1
		{>=}{{\(\ge\)}}1
		{!=}{{\(\neq\)}}1
		{&&}{{\(\wedge\)}}1
		{||}{{\(\vee\)}}1
}

%=========================================================
% Langage LaTeX
%=========================================================
\lstdefinelanguage{latex}{
	morekeywords={
		documentclass,usepackage,input,include,includegraphics,
		begin,end,section,subsection,subsubsection,
		tableofcontents,listoffigures,lstlistoflistings,
		caption,label,ref,vref,textbf,textit,emph,
		DocSetup,newcommand,renewcommand,
		graphicspath,DeclareGraphicsExtensions
	},
	sensitive=true,
	morecomment=[l]{\%},
	morestring=[b]",
	alsoletter={\\}
}

%=========================================================
% Style LaTeX
%=========================================================
\lstdefinestyle{latexdoc}{
	language=latex,
	basicstyle=\ttfamily\small,
	keywordstyle=\color{blue}\bfseries,
	commentstyle=\color{gray},
	stringstyle=\color{red},
	numbers=left,
	numberstyle=\tiny,
	stepnumber=1,
	numbersep=8pt,
	frame=single,
	breaklines=true,
	showstringspaces=false,
	tabsize=2
}

%======================================================================
% PARTIE 8 : Environnement protégé (no copy/paste)
%======================================================================

\lstnewenvironment{lstlistingnocopy}[2]{%
	% #1 : texte remplacé dans le presse-papiers
	% #2 : langage ou style (ex. CSharp, php, mysql, ...)
	\BeginAccSupp{method=escape,ActualText={#1}}%
	\lstset{
		language = #2,
	}%
}{%
	\EndAccSupp{}%
}

%
% USAGE GÉNÉRAL :
%   - Style par défaut neutre :
%       \begin{lstlisting}
%       ...
%       \end{lstlisting}
%
%   - C# (style csharp) :
%       \begin{lstlisting}[style=csharp]
%       ...
%       \end{lstlisting}
%
%   - PHP (style php) :
%       \begin{lstlisting}[style=php]
%       ...
%       \end{lstlisting}
%
%   - JavaScript (style javascript) :
%       \begin{lstlisting}[style=javascript]
%       ...
%       \end{lstlisting}
%
%   - SQL / MySQL (style sql) :
%       \begin{lstlisting}[style=sql]
%       ...
%       \end{lstlisting}
%
%   - Bash (style bashsimple) :
%       \begin{lstlisting}[style=bashsimple]
%       ...
%   \end{lstlisting}
%
%   - Pseudo-code (style pseudocode) :
%       \begin{lstlisting}[style=pseudocode]
%       ...
%   \end{lstlisting}
%
%   - Listing protégé contre le copier/coller :
%       \begin{lstlistingnocopy}{Texte-remplacant-dans-le-presse-papiers}{CSharp}
%       ...
%   \end{lstlistingnocopy}
%
%======================================================================

%----------------------------------------------------------------------
% Tentative de chargement automatique de inc_colors si nécessaire
% On teste une couleur définie dans inc_colors (par ex. dkblue).
%----------------------------------------------------------------------
\makeatletter
\@ifundefined{color@dkblue}{%
	% dkblue non défini -> on suppose que inc_colors n'a pas encore été chargé
	\InputIfFileExists{includes/inc_colors.tex}{}{%
		% Si le fichier n'existe pas, on ne fait rien (fallback silencieux)
	}%
}{}%
\makeatother

%----------------------------------------------------------------------
% Packages nécessaires
%----------------------------------------------------------------------
\usepackage{listings}
\usepackage{listingsutf8} % charge listings + support UTF-8
\usepackage{accsupp}      % pour masquer le texte lors du copy/paste (lstlistingnocopy)

%======================================================================
% PARTIE 1 : Configuration commune
%======================================================================

% Numéros de ligne non copiables (pour certains styles, ex. C#)
\newcommand{\noncopynumber}[1]{%
	\BeginAccSupp{method=escape,ActualText={}}%
	{\scriptsize #1}%
	\EndAccSupp{}%
}

% Réglages globaux pour tous les listings
\lstset{
	basicstyle       = \footnotesize\ttfamily,
	columns          = fullflexible,
	keepspaces       = true,
	showstringspaces = false,
	upquote          = true,
	breaklines       = true,
	breakatwhitespace= true,
	tabsize          = 2,
	frame            = single,
	captionpos       = b,
	inputencoding    = utf8,
	extendedchars    = true,
	% Cartographie commune (accents FR + ponctuation typographique)
	literate         =
		{á}{{\'a}}1 {à}{{\`a}}1 {â}{{\^a}}1 {ä}{{\"a}}1
		{Á}{{\'A}}1 {À}{{\`A}}1 {Â}{{\^A}}1 {Ä}{{\"A}}1
		{é}{{\'e}}1 {è}{{\`e}}1 {ê}{{\^e}}1 {ë}{{\"e}}1
		{É}{{\'E}}1 {È}{{\`E}}1 {Ê}{{\^E}}1 {Ë}{{\"E}}1
		{í}{{\'i}}1 {ì}{{\`i}}1 {î}{{\^i}}1 {ï}{{\"i}}1
		{Í}{{\'I}}1 {Ì}{{\`I}}1 {Î}{{\^I}}1 {Ï}{{\"I}}1
		{ó}{{\'o}}1 {ò}{{\`o}}1 {ô}{{\^o}}1 {ö}{{\"o}}1
		{Ó}{{\'O}}1 {Ò}{{\`O}}1 {Ô}{{\^O}}1 {Ö}{{\"O}}1
		{ú}{{\'u}}1 {ù}{{\`u}}1 {û}{{\^u}}1 {ü}{{\"u}}1
		{Ú}{{\'U}}1 {Ù}{{\`U}}1 {Û}{{\^U}}1 {Ü}{{\"U}}1
		{ç}{{\c{c}}}1 {Ç}{{\c{C}}}1
		{œ}{{\oe}}1 {Œ}{{\OE}}1
		{æ}{{\ae}}1 {Æ}{{\AE}}1
		{’}{{'}}1 {‘}{{`}}1
		{“}{{``}}1 {”}{{''}}1
		{–}{{--}}1 {—}{{---}}1
		{…}{{\ldots}}1
	,
	backgroundcolor  = \color{lightgray},
	rulesepcolor     = \color{gray},
	prebreak         = \raisebox{0ex}[0ex][0ex]{\ensuremath{\hookleftarrow}},
	% Inclusion partielle via rangeprefix/rangesuffix
	rangeprefix        = {//`:~},
	rangesuffix        = {~:`},
	includerangemarker = false,
}

% Style par défaut neutre
\lstdefinestyle{defaultcode}{
	keywordstyle    = \color{dkblue}\bfseries,
	identifierstyle = \color{dkgreen},
	stringstyle     = \color{red},
	commentstyle    = \color{gray}\itshape,
}

\lstset{style=defaultcode}

%======================================================================
% PARTIE 2 : C# (CSharp)
%======================================================================

% Langage C# basé sur [Sharp]C
\lstdefinelanguage{CSharp}{
	language      = [Sharp]C,
	morekeywords  = {
		if, for, while, do, var, get, set,
		string, value, String, ToString, ToInt32, ToBool, toInt
	},
	otherkeywords = {
		\#region, \#endregion, \#define, \#if, \#endif, \#else
	},
}

% Style CSharp complet
\lstdefinestyle{csharp}{
	language        = CSharp,
	numbers         = left,
	numberstyle     = \noncopynumber,
	stepnumber      = 5,
	numbersep       = 5pt,
	frame           = shadowbox,
	keywordstyle    = \color{dkblue}\bfseries,
	commentstyle    = \color{darkgreen},
	stringstyle     = \color{darkred},
	identifierstyle = \color{black},
	% Classes/types mis en avant
	emph            = [1]{
		Application,
		Char, Color, Console, Convert,
		DialogResult,
		Environment, EventArgs,
		Form,
		Object, String,
		SByte, Int16, Int32, Int64,
		Byte, UInt16, UInt32, UInt64,
		Single, Double,
		Keys, KeyPressEventArgs,
		MessageBox, MessageBoxButtons,
		% --- MSTest : classes / attributs importants ---
		Assert,
		TestContext,
		TestClass, TestMethod,
		TestInitialize, TestCleanup,
		ClassInitialize, ClassCleanup,
		DataTestMethod, DataRow,
		ExpectedException,
		TestCategory, Owner, Priority, TestProperty,
		% --- MonoGame / XNA : classes / structs / enums importants ---
		Game, GameTime, GameComponent, DrawableGameComponent,
		GraphicsDevice, GraphicsDeviceManager, GraphicsProfile,
		Viewport, DisplayMode, PresentationParameters,
		ContentManager, ContentReader, ContentWriter,
		SpriteBatch, SpriteFont, SpriteEffects, SpriteSortMode,
		Texture, Texture2D, Texture3D, RenderTarget2D, RenderTargetCube,
		SurfaceFormat, RenderTargetUsage,
		Effect, BasicEffect, AlphaTestEffect,
		BlendState, SamplerState, DepthStencilState, RasterizerState,
		Color, Rectangle, Point,
		Vector2, Vector3, Vector4, Matrix, Quaternion, Plane, Ray,
		MathHelper, Curve, BoundingBox, BoundingSphere, BoundingFrustum,
		GameServiceContainer,
		% --- Input ---
		Keyboard, KeyboardState,
		Mouse, MouseState,
		GamePad, GamePadState, GamePadCapabilities,
		Keys, Buttons, ButtonState, PlayerIndex,
		TouchPanel, TouchCollection, TouchLocation, GestureSample, GestureType,
		% --- Audio / Media ---
		SoundEffect, SoundEffectInstance, SoundState,
		Song, MediaPlayer,
		% --- Timing / misc ---
		TimeSpan,
		% --- Common graphics helpers (si tu utilises les primitives) ---
		VertexDeclaration,
		VertexBuffer, IndexBuffer, DynamicVertexBuffer,
		VertexPositionColor, VertexPositionTexture, VertexPositionColorTexture,
		PrimitiveType, BufferUsage, IndexElementSize
	},
	emphstyle       = [1]\color{class},
}


%======================================================================
% PARTIE 3 : PHP
%======================================================================

% Charger la définition standard de PHP
\lstloadlanguages{PHP}

% Sous-langage PHP "CFPT" : on ajoute nos mots-clés / types
\lstdefinelanguage[CFPT]{PHP}[]{PHP}{%
	% Mots-clés supplémentaires
	morekeywords={declare},
	% Types scalaires / retour qu'on veut colorer comme "types"
	morekeywords=[2]{%
		int, float, string, bool, boolean, void,
		array, iterable, object, mixed, never
	},
}

\lstdefinestyle{php}{%
	language        = [CFPT]PHP,
	% groupe 1 : mots-clés PHP
	keywordstyle    = \color{dkblue}\bfseries,
	% groupe 2 : types
	keywordstyle    = [2]\color{class}\bfseries,
	identifierstyle = \color{dkgreen},
	stringstyle     = \color{red},
	commentstyle    = \color{gray}\itshape,
	emph            = {true, false, null},
	emphstyle       = \color{purple},
	morekeywords    = {
		namespace, use, class, extends,
		public, private, protected, function,
		return, as, new, static, global, var,
		if, else, elseif, while, for, foreach, do,
		switch, case, break, continue, default,
		echo, print, die, exit, instanceof,
		try, catch, finally, throw, implements
	},
	% Délimiteurs spéciaux @@...@@ (par ex. pour marquer du code important)
	moredelim       = **[is][\color{orange}]{@@}{@@},
}

%======================================================================
% PARTIE 4 : JavaScript
%======================================================================

% Définition du langage JavaScript pour listings
\lstdefinelanguage{JavaScript}{
	sensitive        = true,
	% Mots-clés principaux JavaScript
	morekeywords = {
		break, case, catch, continue, debugger, default, delete, do, else, finally,
		for, function, if, in, instanceof, new, return, switch, this, throw, try,
		typeof, var, void, while, with,
		const, let,
		class, extends, super, constructor,
		import, export, from, as,
		get, set,
		async, await, yield
	},
	% Types / "mots-clés" de type (pratique si tu montres un peu de TS ou JSDoc)
	morekeywords = [2]{
		number, string, boolean, any, void, unknown, never, object, Promise
	},
	% Mots-clés / fonctions spécifiques à Jest (groupe 3)
	morekeywords = [3]{
		describe, it, test,
		expect,
		beforeEach, afterEach,
		beforeAll, afterAll,
		jest
	},
	% Commentaires
	morecomment = [l]{//},
	morecomment = [s]{/*}{*/},
	% Chaînes
	morestring  = [b]',
	morestring  = [b]"
}

% Style JavaScript (couleurs cohérentes avec le reste)
\lstdefinestyle{javascript}{%
	language        = JavaScript,
	basicstyle      = \footnotesize\ttfamily,
	% groupe 1 : mots-clés JS
	keywordstyle    = \color{dkblue}\bfseries,
	% groupe 2 : types / notions proches TypeScript
	keywordstyle    = [2]\color{class}\bfseries,
	% groupe 3 : Jest (API de tests)
	keywordstyle    = [3]\color{darkred}\bfseries,
	identifierstyle = \color{dkgreen},
	stringstyle     = \color{red},
	commentstyle    = \color{gray}\itshape,
	emph            = {true, false, null, undefined},
	emphstyle       = \color{purple},
}

%======================================================================
% PARTIE 5 : SQL / MySQL
%======================================================================

\lstdefinelanguage{mysql}{
	morekeywords = {
		% --- requêtes de base ---
		SELECT, FROM, WHERE, INSERT, INTO, VALUES, UPDATE, DELETE,
		CREATE, TABLE, PRIMARY, KEY, AUTO_INCREMENT,
		ALTER, DROP, DATABASE, USE,
		INNER, LEFT, RIGHT, JOIN, ON, AS,
		AND, OR, NOT, NULL, IS, IN, LIKE, BETWEEN,
		ORDER, BY, GROUP, HAVING, DISTINCT, LIMIT, OFFSET,
		FOR,
		% --- transactions / verrous ---
		BEGIN, START, TRANSACTION, COMMIT, ROLLBACK,
		SAVEPOINT, RELEASE, WORK, CHAIN,
		LOCK, UNLOCK,
		% --- niveaux d'isolation / modes d'accès ---
		ISOLATION, LEVEL,
		READ, WRITE, ONLY,
		REPEATABLE, UNCOMMITTED, COMMITTED, SERIALIZABLE,
		AUTOCOMMIT,
		% --- XA transactions ---
		XA, END, PREPARE, RECOVER,
		% --- types ---
		INT, INTEGER, DECIMAL, NUMERIC, FLOAT, DOUBLE,
		CHAR, VARCHAR, TEXT, DATE, DATETIME, TIMESTAMP, TIME,
		% --- contraintes / index ---
		ENGINE, DEFAULT, UNIQUE, INDEX, FOREIGN, REFERENCES,
		% --- gestion des comptes / privilèges ---
		GRANT, REVOKE, ALL, PRIVILEGES,
		USER, ROLE, PASSWORD,
		IDENTIFIED, BY,
		FLUSH, WITH, OPTION, PROXY,
		% --- vues ---
		VIEW, VIEWS,
		% --- procédures / fonctions / triggers / events ---
		PROCEDURE, FUNCTION, TRIGGER, EVENT, ROUTINE,
		INOUT, RETURNS, DETERMINISTIC,
		SQL, SECURITY, DEFINER, INVOKER,
		SET
	},
	sensitive   = false, % mots-clés insensibles à la casse
}

\lstdefinestyle{sql}{
	language        = mysql,
	keywordstyle    = \color{dkblue}\bfseries,
	identifierstyle = \color{dkgreen},
	stringstyle     = \color{red},
	commentstyle    = \color{gray}\itshape,
}

%======================================================================
% PARTIE 6 : Bash / PowerShell (style commun)
%======================================================================

\lstdefinestyle{bashsimple}{
	language        = bash,
	basicstyle      = \footnotesize\ttfamily,
	frame           = single,
	framesep        = 2pt,
	breaklines      = true,
	showstringspaces= false,
	columns         = fullflexible,
	keepspaces      = true,
	upquote         = true,
	tabsize         = 2,
	keywordstyle    = \color{dkblue}\bfseries,
	identifierstyle = \color{black},
	stringstyle     = \color{red},
	commentstyle    = \color{gray}\itshape,
	% IMPORTANT : le tiret fait partie des "lettres"
	alsoletter      = {-},
	% force en noir litéralement
	% IMPORTANT : on conserve la cartographie FR + on ajoute les règles spécifiques
	literate=
		{á}{{\'a}}1 {à}{{\`a}}1 {â}{{\^a}}1 {ä}{{\"a}}1
		{Á}{{\'A}}1 {À}{{\`A}}1 {Â}{{\^A}}1 {Ä}{{\"A}}1
		{é}{{\'e}}1 {è}{{\`e}}1 {ê}{{\^e}}1 {ë}{{\"e}}1
		{É}{{\'E}}1 {È}{{\`E}}1 {Ê}{{\^E}}1 {Ë}{{\"E}}1
		{í}{{\'i}}1 {ì}{{\`i}}1 {î}{{\^i}}1 {ï}{{\"i}}1
		{Í}{{\'I}}1 {Ì}{{\`I}}1 {Î}{{\^I}}1 {Ï}{{\"I}}1
		{ó}{{\'o}}1 {ò}{{\`o}}1 {ô}{{\^o}}1 {ö}{{\"o}}1
		{Ó}{{\'O}}1 {Ò}{{\`O}}1 {Ô}{{\^O}}1 {Ö}{{\"O}}1
		{ú}{{\'u}}1 {ù}{{\`u}}1 {û}{{\^u}}1 {ü}{{\"u}}1
		{Ú}{{\'U}}1 {Ù}{{\`U}}1 {Û}{{\^U}}1 {Ü}{{\"U}}1
		{ç}{{\c{c}}}1 {Ç}{{\c{C}}}1
		{/php/}{{{\color{black}/php/}}}4%
		{/php*}{{{\color{black}/php/}}}4%
		{php.ini}{{{\color{black}php.ini}}}6%
		{.php}{{{\color{black}.php}}}4%
		{wsl.}{{{\color{black}wsl.}}}4
	,
	%--------------------------------------------------------------
	% Commandes bash / sh + cmdlets PowerShell mises en évidence
	% (j'ai retiré 'tar' pour éviter de colorer les extensions .tar)
	%--------------------------------------------------------------
	morekeywords={
		% --- Bash / sh de base ---
		clear,
		cd, ls, cp, mv, rm, mkdir, rmdir, touch,
		chmod, chown, chgrp, umask, usermod,
		sudo, su, passwd,
		apt, apt-get, apt-cache, add-apt-repository,
		a2enmod, a2dismod, a2encond, a2dismod, a2enconf, a2disconf,
		service, systemctl,
		kill, killall, ps, top, htop,
		cat, tac, less, more, head, tail,
		grep, egrep, fgrep, find, xargs,
		echo, printf,
		nano, vi, vim, code,
		ssh, scp, sftp, rsync,
		service,
		zip, unzip, gzip, gunzip,
		df, du, mount, umount,
		uname, hostname, whoami, id, pwd,
		ln,
		php, php8.4,
		mysql, mariadb, mariadb_secure_installation, mysql_secure_installation,
		update-alternatives,
		composer,
		wsl,
		% --- Cmdlets PowerShell fréquents ---
		Get-Help, Get-Command,
		Get-ChildItem, Get-Item, Set-Item, Remove-Item,
		Copy-Item, Move-Item, New-Item, Test-Path,
		Get-Location, Set-Location, Push-Location, Pop-Location,
		Get-Content, Set-Content, Add-Content,
		Get-Process, Stop-Process, Start-Process,
		Get-Service, Start-Service, Stop-Service, Restart-Service,
		Get-Module, Import-Module, Install-Module,
		Invoke-WebRequest, Invoke-Expression,
		Get-History, Clear-History,
		Get-Alias, Set-Alias,
		Get-ExecutionPolicy, Set-ExecutionPolicy,
		Get-EventLog, Write-EventLog
	},
	%--------------------------------------------------------------
	% Options / paramètres mis en avant (bash + PowerShell)
	%--------------------------------------------------------------
	emph={
		% options bash/génériques
		-d,
		-s, status,
		-c,--config,
		-h,--help,
		-v,--version,
		-y,--yes,
		-r,-R,
		-a,--all,
		-l,--list,
		-f,--force,
		--update,
		-i,--ini,
		--install,install,
		--export, --terminate,
		--import,
		-name,
		restart,
		--unregister,
		--set-default,
		--list,
		--verbose,
		-G,
		-u, --user,
		-p, --password,
		update, full-upgrade,
		% options courantes PowerShell
		-Path,
		-Recurse,
		-Filter,
		-Include,
		-Exclude,
		-Force,
		-Name,
		-ErrorAction,
		-WhatIf,
		-Confirm,
		-File,
		-Directory,
		-Value,
		-Destination,
		-Source,
		-Credential
	},
	emphstyle=\color{purple}\bfseries,
}

%======================================================================
% Style listings : Pseudo-code (Flowgorithm / ATALGO)
%======================================================================

\lstdefinestyle{pseudocode}{
	language        = {},
	basicstyle      = \footnotesize\ttfamily,
	showstringspaces=false,
	columns         = fullflexible,
	keepspaces      = true,
	upquote         = true,
	% Commentaires (si tu en mets dans tes extraits)
	morecomment=[l]{//},
	morecomment=[s]{/*}{*/},
	commentstyle    = \color{gray}\itshape,
	% Chaînes
	morestring=[b]",
	stringstyle     = \color{darkred},
	% -------------------------
	% Groupe 1 : Structure / contrôle
	% -------------------------
	keywordstyle    = \color{dkblue}\bfseries,
	morekeywords    = {
		ALGORITHME, DEBUT_ALGORITHME, FIN_ALGORITHME,
		Fonction, Principale, DEBUT, FIN, Fin,
		SI, Si, ALORS, Alors, SINON, Sinon, Autre, DEBUT_SI, FIN_SI, DEBUT_SINON, FIN_SINON,
		TANT_QUE, Tant, Que, DEBUT_TANT_QUE, FIN_TANT_QUE,
		POUR, Pour, DEBUT_POUR, FIN_POUR, ALLANT_DE, PAR_PAS_DE, A, a,
		FAIRE, Faire,
		CAS, SUIVANT, DEBUT_SUIVANT, FIN_SUIVANT
	},
	% -------------------------
	% Groupe 2 : I/O et actions “verbes” (couleur différente)
	% -------------------------
	classoffset=1,
	keywordstyle=\color{darkgreen}\bfseries,
	morekeywords={
		AFFICHER, Afficher,
		LIRE, Lire,
		ATTENDRE, Attendre,
		Saisie, SAISIE,
		Appeler, APPELER,
		Initialiser, INITIALISER
	},
	classoffset=0,
	% -------------------------
	% Groupe 3 : Déclarations / affectations (encore une autre couleur)
	% -------------------------
	classoffset=2,
	keywordstyle=\color{dkblue}\bfseries,
	morekeywords={
		VARIABLES, Variables,
		Déclarer, Declarer, DECLARER,
		Variable,
		Type, TYPE, de,
		Assigner, ASSIGNER,
		Expression, EXPRESSION,
		Paramètre, Parametre, PARAMETRE,
		Retour, RETOUR,
		Retourner, Retourner, RETOURNER,
		Retouner, RETOUNER, % (typo présent dans le PDF)
		Prend, PREND, PREND_POUR_VALEUR
	},
	classoffset=0,
	% -------------------------
	% Types fréquents (mis en “emphase”)
	% -------------------------
	emph={Entier,Réel,Reel,Chaîne,Chaine,Booléen,Booleen,Tableau,Array},
	emphstyle=\color{darkred}\bfseries,
	% Fonctions usuelles vues dans tes pseudo-codes
	emph={size,len,char,Random,True,False},
	emphstyle=\color{darkgreen},
	% Bonus : rendre quelques opérateurs plus “propres” (optionnel)
	literate=
		{á}{{\'a}}1 {à}{{\`a}}1 {â}{{\^a}}1 {ä}{{\"a}}1
		{Á}{{\'A}}1 {À}{{\`A}}1 {Â}{{\^A}}1 {Ä}{{\"A}}1
		{é}{{\'e}}1 {è}{{\`e}}1 {ê}{{\^e}}1 {ë}{{\"e}}1
		{É}{{\'E}}1 {È}{{\`E}}1 {Ê}{{\^E}}1 {Ë}{{\"E}}1
		{í}{{\'i}}1 {ì}{{\`i}}1 {î}{{\^i}}1 {ï}{{\"i}}1
		{Í}{{\'I}}1 {Ì}{{\`I}}1 {Î}{{\^I}}1 {Ï}{{\"I}}1
		{ó}{{\'o}}1 {ò}{{\`o}}1 {ô}{{\^o}}1 {ö}{{\"o}}1
		{Ó}{{\'O}}1 {Ò}{{\`O}}1 {Ô}{{\^O}}1 {Ö}{{\"O}}1
		{ú}{{\'u}}1 {ù}{{\`u}}1 {û}{{\^u}}1 {ü}{{\"u}}1
		{Ú}{{\'U}}1 {Ù}{{\`U}}1 {Û}{{\^U}}1 {Ü}{{\"U}}1
		{ç}{{\c{c}}}1 {Ç}{{\c{C}}}1
		{<=}{{\(\le\)}}1
		{>=}{{\(\ge\)}}1
		{!=}{{\(\neq\)}}1
		{&&}{{\(\wedge\)}}1
		{||}{{\(\vee\)}}1
}

%=========================================================
% Langage LaTeX
%=========================================================
\lstdefinelanguage{latex}{
	morekeywords={
		documentclass,usepackage,input,include,includegraphics,
		begin,end,section,subsection,subsubsection,
		tableofcontents,listoffigures,lstlistoflistings,
		caption,label,ref,vref,textbf,textit,emph,
		DocSetup,newcommand,renewcommand,
		graphicspath,DeclareGraphicsExtensions
	},
	sensitive=true,
	morecomment=[l]{\%},
	morestring=[b]",
	alsoletter={\\}
}

%=========================================================
% Style LaTeX
%=========================================================
\lstdefinestyle{latexdoc}{
	language=latex,
	basicstyle=\ttfamily\small,
	keywordstyle=\color{blue}\bfseries,
	commentstyle=\color{gray},
	stringstyle=\color{red},
	numbers=left,
	numberstyle=\tiny,
	stepnumber=1,
	numbersep=8pt,
	frame=single,
	breaklines=true,
	showstringspaces=false,
	tabsize=2
}

%======================================================================
% PARTIE 8 : Environnement protégé (no copy/paste)
%======================================================================

\lstnewenvironment{lstlistingnocopy}[2]{%
	% #1 : texte remplacé dans le presse-papiers
	% #2 : langage ou style (ex. CSharp, php, mysql, ...)
	\BeginAccSupp{method=escape,ActualText={#1}}%
	\lstset{
		language = #2,
	}%
}{%
	\EndAccSupp{}%
}

%-------------------------------------------------------------------------------------
% Métadonnées PDF et style des liens (hyperref)
%-------------------------------------------------------------------------------------
\hypersetup{
	% Titre du document dans les propriétés PDF
	pdftitle={\DocTitle{} : \DocSubtitle{}},
	% Auteur du document (nom + email) dans les propriétés PDF
	pdfauthor={\DocAuthor{} -- \DocEmail{}},
	% Sujet / description courte (domaine + année scolaire)
	pdfsubject={\DocDomaine{} -- \DocYear{}},
	% Mots-clés pour l’indexation du PDF (recherche, catalogage, etc.)
	pdfkeywords={\DocKeywords},
	% Informations techniques sur la génération du PDF (producteur)
	pdfproducer={MiKTeX + pdfTeX}, % facultatif : rempli automatiquement si omis
	% Informations techniques sur la génération du PDF (créateur)
	pdfcreator={LaTeX with hyperref}, % facultatif : idem
	% Style des liens : true = liens colorés, false = liens encadrés
	colorlinks=true,
	% Couleur des liens internes (table des matières, renvois, etc.)
	linkcolor=blue,
	% Couleur des liens externes (URL, adresses web)
	urlcolor=blue
}
